% This file was converted to LaTeX by Writer2LaTeX ver. 1.9.9
% see http://writer2latex.sourceforge.net for more info
\documentclass{article}
\usepackage{calc,amsmath,amssymb,amsfonts}
\usepackage[T1]{fontenc}
\usepackage[italian]{babel}
\usepackage[style=numeric,backend=biber]{biblatex}
\author{HP}
\date{2026-07-26}
\begin{document}
\section{Prologo}
Ci sono ricordi che il tempo sfuma lentamente, fino a renderli quasi irriconoscibili.

Altri, invece, rimangono sorprendentemente nitidi. Non sapresti dire che giorno fosse, né ricordare con precisione il
mese o l'ora. Eppure, se chiudi gli occhi, riesci ancora a rivedere la stanza, la luce che entrava dalla finestra e
perfino la sensazione provata in quell'istante.

Il mio primo incontro con un videogioco appartiene a questa seconda categoria.

Era il 1982.

Frequentavo la scuola media e un pomeriggio andai a trovare un mio compagno di classe, Emanuele. Era una visita come
tante altre. Avremmo trascorso qualche ora insieme, parlando di scuola, giocando o semplicemente facendo quello che
fanno due ragazzi della stessa età.

Non avevo la minima idea che quel pomeriggio sarebbe rimasto impresso nella mia memoria per tutta la vita.

Entrai in casa, salutai i suoi genitori e lo seguii nella sua cameretta.

Ricordo una stanza ordinata, forse più di quanto lo fosse la mia. La scrivania era appoggiata contro una parete e,
sopra, c'era un televisore. Niente lasciava immaginare che da lì a pochi minuti avrei visto qualcosa che avrebbe
cambiato per sempre il mio modo di guardare un apparecchio televisivo.

Per me, fino a quel momento, la televisione aveva avuto un solo scopo: trasmettere immagini.

Cartoni animati.

Telefilm.

Film.

Telegiornali.

La si accendeva per guardare ciò che qualcun altro aveva deciso di mandare in onda.

Quello che vidi quel pomeriggio era completamente diverso.

Sul televisore comparvero due barre bianche e un piccolo quadrato luminoso che si muoveva da una parte all'altra dello
schermo.

Era Pong.

Oggi può sembrare difficile comprenderlo. Abituati a immagini fotorealistiche, mondi tridimensionali e console capaci di
elaborare miliardi di operazioni al secondo, quel gioco potrebbe apparire quasi banale.

Ma nel 1982 non era così.

Per me non era soltanto un gioco.

Era una domanda.

Guardavo quella piccola pallina rimbalzare contro le racchette e cercavo di capire.

Come faceva?

Come poteva un televisore sapere dove si trovava quella pallina?

Come faceva a cambiarne direzione quando colpiva una racchetta?

Come decideva quando assegnare un punto?

Più osservavo lo schermo, meno mi interessava vincere.

Volevo capire.

Probabilmente Emanuele se ne accorse. Mentre lui giocava, io continuavo a osservare il televisore con lo stesso stupore
con cui si guarda qualcosa che non si riesce ancora a spiegare.

Quel pomeriggio finì come tutti i pomeriggi.

Salutai.

Uscii di casa.

Percorsi la strada verso casa senza pensarci troppo.

O almeno così mi sembrava.

Perché, in realtà, una domanda aveva già iniziato a seguirmi.

Continuava a ripresentarsi, ostinata, senza chiedere il permesso.

Come fa?

Negli anni successivi quella domanda non scomparve.

Anzi, diventò sempre più insistente.

Ogni volta che vedevo un videogioco, ogni volta che osservavo un computer, ogni volta che sentivo parlare di
programmazione, tornava a bussare nella mia mente.

Come fa?

Non immaginavo che la risposta sarebbe arrivata l'anno successivo, grazie a un professore che non avevo mai incontrato e
a un computer che portava un numero apparentemente insignificante.

Sessantaquattro.


\bigskip
\end{document}
